%% -*- coding: utf-8 -*-
\documentclass[12pt,pagesize,paper=192mm:108mm,landscape]{scrbook} 
%1920x1080 1280x720
\areaset[current]{192mm}{108mm}
\usepackage{calc}
\usepackage[T2A]{fontenc}
\usepackage[utf8]{inputenc}
\usepackage[english,russian]{babel}
\usepackage{microtype}
\usepackage{misccorr}
\usepackage{cmap}
%\usepackage[unicode=true]{hyperref}
\usepackage{graphicx}
\usepackage{amssymb}
\usepackage{amsmath}
%\usepackage{srcltx}
\usepackage{textcomp}
\usepackage{xspace}
%научные символы и смайлики \smiley \frownie
\usepackage{wasysym}
\usepackage{ccicons}
\begin{document}
\begin{titlepage}
  \vspace*{-0.5em}
  \begin{center}    
    \hspace*{3em}
    \begin{minipage}[t]{1.5em}
      \includegraphics[width=\textwidth]{../BINP-logo}
    \end{minipage}\hfill
    \begin{minipage}{0.115\linewidth}
    \includegraphics[width=\textwidth]{../NSU-logo}
    \end{minipage}
    \hfill
    \hspace*{4.5em}

    Кафедра теоретической физики физического факультета НГУ
    \medskip

    \Large
    Профессор Грабовский А.\,В.
    \smallskip

    \huge
    \textbf{Стандартная модель}
    \smallskip

    \Large
    Лекция № 3
    \vfill

    \normalsize
    \begin{minipage}{0.85\linewidth}
      Лагранжиан бозона Хиггса, симметрия вакуумного состояния
      относительно группы электрического заряда $U(1)_Q$. Спонтанное
      нарушение $U(1)_ {Y_W} \times SU(2)_{I_W}$ до $U(1)_Q$. Механизм
      Хиггса. Количество степеней свободы у полей до и после
      спонтанного нарушения симметрии. Масса бозона Хиггса. Лагранжиан
      взаимодействия бозона Хиггса и фермионов. Лагранжиан
      взаимодействия бозона Хиггса и векторных бозонов. Вершины
      взаимодействия трех и четырех калибровочных бозонов. Охранная
      (custodial) симметрия потенциала поля Хиггса. Параметр $\rho$,
      определяющий отношений силы заряженного и нейтрального токов при
      низких энергиях.
    \end{minipage}
    \vfill

    \normalsize \ccbysa\hspace{0.5em}  Новосибирск 2024
  \end{center}
\end{titlepage}
\end{document}
